\documentclass[11pt,letterpaper]{article}
\usepackage[utf8]{inputenc}
\usepackage{geometry}
\usepackage{amsmath}
\usepackage{setspace}
\usepackage{url}
\usepackage{enumitem}
\usepackage{hyperref}

\geometry{left=1.0in,right=1.0in,top=1.0in,bottom=1.0in}
\onehalfspacing

\setlength{\parindent}{0pt}
\setlength{\parskip}{1.0em plus 0.2em minus 0.2em}

\begin{document}

\thispagestyle{empty}

\begin{center}
    \textbf{\large IN THE UNITED STATES PATENT AND TRADEMARK OFFICE}
\end{center}

\vspace{0.3cm}
\noindent\begin{tabular}{p{2.3in}p{4.0in}}
\textbf{Applicant / Inventor:} & Reinaldo Javier Aguiar Marcano \\
\textbf{Application Number:} & 19/640,945 \\
\textbf{Filing Date:} & April 7, 2026 \\
\textbf{Confirmation Number:} & 7462 \\
\textbf{Art Unit / Examiner:} & TC 2100 / ANDREA BENNETT \\
\textbf{Office of Petitions:} & Charles Steven Brantley, Attorney Advisor \\
\textbf{Customer Number:} & 226458 \\
\textbf{Entity Status:} & SMALL ENTITY (37 CFR \S{} 1.27) \\
\textbf{Title of Invention:} & METHOD FOR ASYNCHRONOUS SECURE COMMUNICATION \\
& WITH POST-HOC KEY DERIVATION \\
\end{tabular}

\vspace{0.4cm}
\begin{center}
    \textbf{\large APPLICANT RESPONSE TO NOTICE OF INCOMPLETE REPLY \\
    (37 CFR \S{} 1.53(b))} \\
    \textbf{\normalsize (SUBMISSION OF REPLACEMENT CLAIMS, REPLACEMENT ABSTRACT, \\
    FEE ALLOCATION, AND CONDITIONAL SUSPENSION)}
\end{center}

\noindent\textbf{Mail Stop Missing Parts} \\
Commissioner for Patents \\
P.O. Box 1450 \\
Alexandria, VA 22313-1450

\vspace{0.2cm}
Sir:

In responsive reply to the \textit{Notice of Incomplete Reply} mailed October 1, 2026, and the \textit{Decision on Petition} (PETDEC) mailed October 8, 2026, the pro se Applicant, Reinaldo Javier Aguiar Marcano, respectfully submits this complete formalities response.

\section*{I. STATEMENT OF REAFFIRMED ENTITY STATUS (37 CFR \S{} 1.27)}

Applicant hereby formally reaffirms entitlement to \textbf{Small Entity status} under 37 CFR \S{} 1.27 and 35 U.S.C. \S{} 123. Applicant's verified Statement of Small Entity Status has been officially recognized on the docket by the Office. All statutory fees are computed strictly under the Small Entity schedule.

\section*{II. CURING OF FORMALITY DEFECTS \& COMPLETE ITEMIZED SUBMISSIONS}

To bring the application into 100\% compliance with USPTO pre-examination requirements, Applicant concurrently submits the following standalone statutory papers in this electronic submission envelope:

\begin{enumerate}[leftmargin=*]
    \item \textbf{Standalone Replacement Claims (37 CFR \S{} 1.75(h), \S{} 1.121, \S{} 1.126) [Patent Center Description: ``Claims'']:} \\
    Applicant submits clean replacement Claims 1--6 commencing on a separate sheet, formatted in strict double line spacing pursuant to 37 CFR \S{} 1.52(b)(2)(i), consecutively numbered in Arabic numerals without duplication, under the pure statutory heading ``CLAIMS'' and introductory phrase ``What is claimed is:''.

    \item \textbf{Standalone Replacement Abstract (37 CFR \S{} 1.72(b), \S{} 1.121) [Patent Center Description: ``Abstract'']:} \\
    Applicant submits a clean replacement Abstract commencing on a separate sheet, formatted in strict double line spacing, under 150 words in length (115 words), under the pure statutory heading ``ABSTRACT OF THE DISCLOSURE''.

    \item \textbf{Acknowledged Inventor's Declaration (37 CFR \S{} 1.63) [Patent Center Description: ``Oath or Declaration filed'']:} \\
    Applicant notes that an executed Inventor's Oath or Declaration under 37 CFR \S{} 1.63 was filed at initial filing and acknowledged in the Image File Wrapper.

    \item \textbf{Conditional Petition for Extension of Time (37 CFR \S{} 1.136(a)) [Patent Center Description: ``Extension of Time'']:} \\
    Applicant submits a Conditional Petition for Extension of Time under 37 CFR \S{} 1.136(a). The initial two-month reply period set in the August 18, 2026 Notice expires on Sunday, October 18, 2026, which extends automatically to \textbf{Monday, October 19, 2026} pursuant to 35 U.S.C. \S{} 21 and 37 CFR \S{} 1.7(a). Because this reply is filed on October 10, 2026, it is fully timely within the initial statutory window, requiring \$0.00 extension fees today.

    \item \textbf{Conditional Request for Six-Month Suspension (37 CFR \S{} 1.103(a)) [Patent Center Description: ``Letter Requesting Suspension of Action'']:} \\
    Applicant submits a Conditional Request for Suspension of Action under 37 CFR \S{} 1.103(a) (Fee Code 2830, \$60.00), requesting that prosecution be suspended for six (6) months upon completion of pre-examination processing, due to ongoing federal civil rights and constitutional litigation in the United States District Court for the Southern District of Texas, Houston Division, Case No. 4:26-cv-08794 (Judge Kenneth M. Hoyt; Cause: 28:1983 Civil Rights, Nature of Suit: 440).
\end{enumerate}

\section*{III. STATUTORY FEE RECONCILIATION, REQUEST FOR DEPOSIT ALLOCATION (37 CFR \S{} 1.25(b)), AND MACRO PORTFOLIO EQUITY}

\subsection*{A. Macro Portfolio Fee Equity Context}
Applicant respectfully directs the Office of Patent Application Processing and the Office of Finance to the overarching financial realities of this prosecution. Across the 102 nonprovisional patent applications comprising Applicant's master portfolio, the USPTO Office of Finance is currently holding approximately \$35,000.00 to \$42,000.00 USD (with \$32,758.00 USD already empirically verified and credited across 79 examined notices in the Image File Wrappers) in unapplied funds on deposit, without interest and without refund.

Viewing this individual application through an isolated, single-docket lens distorts the financial record by suggesting that Applicant has failed to remit fees. In truth, Applicant has already remitted substantial capital that remains on deposit with the agency. Fundamental administrative equity and MPEP \S{} 509 require that unapplied deposit credits be properly accounted for and applied to maintain procedural timeliness, rather than assessing progressive extension fees that compound to over \$400,000 across the portfolio without prior fee reconciliation.

\subsection*{B. Respectful Request for Allocation of Unapplied Deposit Funds (37 CFR \S{} 1.25(b) \& MPEP \S{} 509)}
Official Office records (specifically the October 1, 2026 Notice of Incomplete Reply and Fee Worksheet WFEE) confirm that the Office of Finance currently holds \textbf{\$450.00 in unapplied funds on deposit} under this application serial number.

Concurrently herewith, Applicant remits the \textbf{\$60.00 processing fee} under 37 CFR \S{} 1.17(i) (Fee Code 2830) via electronic fee payment in Patent Center to satisfy the statutory requirement for the accompanying Request for Suspension under 37 CFR \S{} 1.103(a).

Pursuant to 37 CFR \S{} 1.25(b), MPEP \S{} 509, and MPEP \S{} 710.02(e), Applicant hereby respectfully requests that the Office allocate said \$450.00 unapplied funds on deposit in the following order of procedural priority:
\begin{enumerate}[leftmargin=*]
    \item \textbf{PRIORITY 1 (Pre-Exam Base Fee Delta Reduction):} Because this response is filed on October 10, 2026, well prior to the October 19, 2026 deadline, no extension fee is currently due. Applicant respectfully requests that the entire available deposit credit of \textbf{\$450.00} be applied directly toward the outstanding pre-examination base fee balance under 37 CFR \S{} 1.16, reducing the net base fee deficiency from \$520.00 to \textbf{\$70.00}. Because this response is timely filed and substantially funded, Applicant respectfully requests that the Office issue a formal Notice of Fee Deficiency under 37 CFR \S{} 1.53(f) and MPEP \S{} 506 for the remaining base balance of \$70.00, or hold collection in abeyance pursuant to the concurrent Request for Suspension under 37 CFR \S{} 1.103(a);
    \item \textbf{PRIORITY 2 (Suspension Processing):} Concurrently herewith, Applicant remits the \$60.00 processing fee (Fee Code 2830) by electronic fee payment in Patent Center to satisfy the statutory requirement for the accompanying Request for Suspension under 37 CFR \S{} 1.103(a).
\end{enumerate}

\section*{IV. CONCLUSION}

All formality requirements having been satisfied, Applicant respectfully requests that the Office of Patent Application Processing (OPAP) advance the application to Technology Center 2100 for examination.

\vspace{0.6cm}
\begin{singlespace}
\noindent Respectfully submitted, \par\vspace{0.8cm}
\underline{\textbf{/Reinaldo Javier Aguiar Marcano/}} \\
\textbf{Reinaldo Javier Aguiar Marcano} \\
Inventor and Pro Se Applicant \\
Mailing Address: 2302 Britton Ridge Drive, Katy, Texas 77494 \\
Telephone: (202) 599-3979 \\
Email: \texttt{rj.aguiar@gmail.com} (USPTO Account) \\
Email: \texttt{rj.aguiar.m@gmail.com} \\
Email: \texttt{reinaldo.aguiar.marcano@isbn9781971481296.com} \\
Date: October 10, 2026
\end{singlespace}

\end{document}
